\FigureLayoutDeclare{img/2001/new_curriculum_science/q11_fig1.png}{15.0cm}{0bp 0bp 0bp 0bp}
\FigureLayoutDeclare{img/2001/new_curriculum_science/q20_fig2.png}{6.5cm}{0bp 0bp 0bp 0bp}

\chapter{2001年新课标卷（理）}
\section{单选题}
\begin{problem}
函数 \(y=3\sin\left(\frac{x}{2}+\frac{\pi}{3}\right)\) 的周期，振幅依次是（\quad）

\choices
    {\(4\pi,3\)}
    {\(4\pi\)，\(-3\)}
    {\(\pi\)，\(3\)}
    {\(\pi\)，\(-3\)}
\end{problem}

\begin{answer}
A
\end{answer}

\begin{solution}
正弦函数 \(y=A\sin(\omega x+\varphi)\) 的振幅为 \(\lvert A\rvert\)，周期为 \(\dfrac{2\pi}{\lvert\omega\rvert}\). 题中 \(A=3\)，\(\omega=\dfrac12\)，所以振幅为 \(\lvert3\rvert=3\)，周期为
\[
T=\frac{2\pi}{\left\lvert\frac12\right\rvert}=4\pi
\]
因此选 A.
\end{solution}

\begin{problem}
若 \(S_n\) 是数列 \(\{a_n\}\) 的前 \(n\) 项和，且 \(S_n=n^2\)，则 \(\{a_n\}\) 是（\quad）

\choices
    {等比数列，但不是等差数列}
    {等差数列，但不是等比数列}
    {等差数列，而且也是等比数列}
    {既非等比数列又非等差数列}
\end{problem}

\begin{answer}
B
\end{answer}

\begin{solution}
由 \(S_1=1^2=1\)，得 \(a_1=S_1=1\). 当 \(n\geqslant2\) 时，
\[
a_n=S_n-S_{n-1}=n^2-(n-1)^2=2n-1
\]
所以 \(a_n=2n-1\)，数列为 \(1\)，\(3\)，\(5\)，\(\ldots\)，且
\[
a_{n+1}-a_n=\bigl(2(n+1)-1\bigr)-(2n-1)=2
\]
故它是等差数列. 另一方面，\(\dfrac{a_2}{a_1}=3\)，而 \(\dfrac{a_3}{a_2}=\dfrac53\)，相邻两项的比不恒定，故它不是等比数列. 因此选 B.
\end{solution}

\begin{problem}
过点 \(A(1,-1)\)，\(B(-1,1)\) 且圆心在直线 \(x+y-2=0\) 上的圆的方程是（\quad）

\choices
    {\((x-3)^2+(y+1)^2=4\)}
    {\((x+3)^2+(y-1)^2=4\)}
    {\((x-1)^2+(y-1)^2=4\)}
    {\((x+1)^2+(y+1)^2=4\)}
\end{problem}

\begin{answer}
C
\end{answer}

\begin{solution}
设圆心为 \(M(x_0,y_0)\). 圆经过 \(A\)，\(B\)，所以 \(M\) 在弦 \(AB\) 的垂直平分线上. 线段 \(AB\) 的中点为 \(O(0,0)\)，直线 \(AB\) 的斜率为 \(-1\)，故其垂直平分线为 \(y=x\). 又 \(M\) 在 \(x+y-2=0\) 上，于是
\(y_0=x_0\)，\(x_0+y_0=2\)
从而 \(M(1,1)\). 半径的平方为
\[
r^2=MA^2=(1-1)^2+(1-(-1))^2=4
\]
因此圆的方程为
\[
(x-1)^2+(y-1)^2=4
\]
故选 C.
\end{solution}

\begin{problem}
若定义在区间 \((-1,0)\) 内的函数 \(f(x)=\log_{2a}(x+1)\) 满足 \(f(x)>0\)，则 \(a\) 的取值范围是（\quad）

\choices
    {\(\left(0,\frac{1}{2}\right)\)}
    {\(\left(0,\frac{1}{2}\right]\)}
    {\(\left(\frac{1}{2},+\infty\right)\)}
    {\((0,+\infty)\)}
\end{problem}

\begin{answer}
A
\end{answer}

\begin{solution}
当 \(x\in(-1,0)\) 时，\(0<x+1<1\). 对于 \(0<u<1\)，当且仅当对数的底满足 \(0<2a<1\) 时，\(\log_{2a}u>0\). 因而
\[
0<2a<1\quad\Longleftrightarrow\quad 0<a<\frac12
\]
故选 A.
\end{solution}

\begin{problem}
若向量 \(\bs{a}=(1,1)\)，\(\bs{b}=(1,-1)\)，\(\bs{c}=(-1,2)\)，则 \(\bs{c}=\)（\quad）

\choices
    {\(-\frac{1}{2}\bs{a}+\frac{3}{2}\bs{b}\)}
    {\(\frac{1}{2}\bs{a}-\frac{3}{2}\bs{b}\)}
    {\(\frac{3}{2}\bs{a}-\frac{1}{2}\bs{b}\)}
    {\(-\frac{3}{2}\bs{a}+\frac{1}{2}\bs{b}\)}
\end{problem}

\begin{answer}
B
\end{answer}

\begin{solution}
设 \(\bs{c}=m\bs{a}+n\bs{b}\). 比较两个坐标，得到
\(m+n=-1\)，\(m-n=2\).
两式相加得 \(2m=1\)，所以 \(m=\dfrac12\)；两式相减得 \(2n=-3\)，所以 \(n=-\dfrac32\). 因而
\[
\bs{c}=\frac12\bs{a}-\frac32\bs{b}
\]
故选 B.
\end{solution}

\begin{problem}
设 \(A\)，\(B\) 是 \(x\) 轴上的两点，点 \(P\) 的横坐标为 \(2\) 且 \(\lvert PA\rvert=\lvert PB\rvert\). 若直线 \(PA\) 的方程为 \(x-y+1=0\)，则直线 \(PB\) 的方程是（\quad）

\choices
    {\(x+y-5=0\)}
    {\(2x-y-1=0\)}
    {\(2y-x-4=0\)}
    {\(2x+y-7=0\)}
\end{problem}

\begin{answer}
A
\end{answer}

\begin{solution}
直线 \(PA\) 与 \(x\) 轴的交点为 \(A(-1,0)\). 因为 \(P\) 的横坐标为 \(2\)，且 \(P\in PA\)，代入直线方程得 \(P(2,3)\). 于是
\[
PA^2=(2-(-1))^2+(3-0)^2=18
\]
设 \(B(t,0)\). 由 \(\lvert PB\rvert=\lvert PA\rvert\)，得
\((t-2)^2+3^2=18\)，\(\Longrightarrow\)，\((t-2)^2=9\)
所以 \(t=-1\) 或 \(t=5\). 由于 \(A\)，\(B\) 是两点，\(B\ne A\)，故 \(B(5,0)\). 直线 \(PB\) 的斜率为 \(-1\)，其方程为
\[
y-3=-(x-2)
\]
即 \(x+y-5=0\). 故选 A.
\end{solution}

\begin{problem}
若 \(0<\alpha<\beta<\frac{\pi}{4}\)，\(\sin\alpha+\cos\alpha=a\)，\(\sin\beta+\cos\beta=b\)，则（\quad）

\choices
    {\(a<b\)}
    {\(a>b\)}
    {\(ab<1\)}
    {\(ab>2\)}
\end{problem}

\begin{answer}
A
\end{answer}

\begin{solution}
令 \(g(t)=\sin t+\cos t\). 当 \(0<t<\dfrac\pi4\) 时，\(\sin t<\cos t\)，所以
\[
g'(t)=\cos t-\sin t>0
\]
因此 \(g(t)\) 在 \(\left(0,\dfrac\pi4\right)\) 上严格递增. 由 \(\alpha<\beta\)，有
\[
a=g(\alpha)<g(\beta)=b
\]
故选 A.
\end{solution}

\begin{problem}
函数 \(y=1+3x-x^3\) 有（\quad）

\choices
    {极小值 \(-1\)，极大值 \(1\)}
    {极小值 \(-2\)，极大值 \(3\)}
    {极小值 \(-2\)，极大值 \(2\)}
    {极小值 \(-1\)，极大值 \(3\)}
\end{problem}

\begin{answer}
D
\end{answer}

\begin{solution}
函数的导数为
\[
y'=3-3x^2=3(1-x)(1+x)
\]
当 \(x<-1\) 或 \(x>1\) 时，\(y'<0\)；当 \(-1<x<1\) 时，\(y'>0\). 所以函数在 \((-\infty,-1)\) 上递减，在 \((-1,1)\) 上递增，在 \((1,+\infty)\) 上递减. 因此 \(x=-1\) 处取得极小值，\(x=1\) 处取得极大值，且
\(y(-1)=1-3+1=-1\)，\(y(1)=1+3-1=3\).
故极小值为 \(-1\)，极大值为 \(3\)，选 D.
\end{solution}

\begin{problem}
某赛季足球比赛的计分规则是：胜一场，得 \(3\text{ 分}\)；平一场，得 \(1\text{ 分}\)；负一场，得 \(0\text{ 分}\). 一球队打完 \(15\text{ 场}\)，积 \(33\text{ 分}\). 若不考虑顺序，该队胜，负，平的情况共有（\quad）

\choices
    {\(3\text{ 种}\)}
    {\(4\text{ 种}\)}
    {\(5\text{ 种}\)}
    {\(6\text{ 种}\)}
\end{problem}

\begin{answer}
A
\end{answer}

\begin{solution}
设胜、负、平的场数分别为 \(w\)，\(l\)，\(d\)，则
\(w+l+d=15\)，\(3w+d=33\).
由第二式得 \(d=33-3w\)，代入第一式得 \(l=2w-18\). 因为 \(w\)，\(l\)，\(d\) 都是非负整数，所以
\(2w-18\geqslant0\)，\(33-3w\geqslant0\)
即 \(9\leqslant w\leqslant11\). 逐一取 \(w=9\)，\(10\)，\(11\)，得到
\((w,l,d)=(9,0,6)\)，\((10,2,3)\)，\((11,4,0)\).
因此共有 \(3\) 种情况，选 A.
\end{solution}

\begin{problem}
设坐标原点为 \(O\)，抛物线 \(y^2=2x\) 与过焦点的直线交于 \(A\)，\(B\) 两点，则 \(\overrightarrow{OA}\cdot\overrightarrow{OB}=\)（\quad）

\choices
    {\(\frac{3}{4}\)}
    {\(-\frac{3}{4}\)}
    {\(3\)}
    {\(-3\)}
\end{problem}

\begin{answer}
B
\end{answer}

\begin{solution}
将抛物线写成 \(y^2=4px\)，得 \(p=\dfrac12\)，所以焦点为 \(F\left(\dfrac12,0\right)\).

先考虑过焦点且不是竖直直线的情形. 由于直线与抛物线交于两个不同的点，它不是 \(x\) 轴，因此可写成
\[
x=my+\frac12
\]
设交点为 \(A(x_1,y_1)\)，\(B(x_2,y_2)\). 联立直线与抛物线，得
\[
y^2=2my+1
\]
即
\[
y^2-2my-1=0
\]
由韦达定理，
\(y_1+y_2=2m\)，\(y_1y_2=-1\).
又 \(x_i=my_i+\dfrac12\)，所以
\[
x_1x_2=m^2y_1y_2+\frac m2(y_1+y_2)+\frac14
=-m^2+m^2+\frac14=\frac14
\]
从而
\[
\overrightarrow{OA}\cdot\overrightarrow{OB}=x_1x_2+y_1y_2=\frac14-1=-\frac34
\]

若过焦点的直线是竖直直线 \(x=\dfrac12\)，则与抛物线联立得 \(y^2=1\). 这时两个交点的坐标为
\(\left(\dfrac12,1\right)\) 和 \(\left(\dfrac12,-1\right)\)，仍有
\[
\overrightarrow{OA}\cdot\overrightarrow{OB}=\frac14-1=-\frac34
\]
因此选 B.
\end{solution}

\begin{problem}
一间民房的屋顶有如图三种不同的盖法，分别是：

\begin{circlelist}
\item 单向倾斜；
\item 双向倾斜；
\item 四向倾斜.
\end{circlelist}

记三种盖法屋顶面积分别为 \(P_1\)，\(P_2\)，\(P_3\). 若屋顶斜面与水平面所成的角都是 \(\alpha\)，则（\quad）

\begingroup
\leftskip=0pt\relax
\rightskip=0pt\relax
\bitmapfigure[max width=0.74\linewidth]{img/2001/new_curriculum_science/q11_fig1.png}
\endgroup

\choices
    {\(P_3>P_2>P_1\)}
    {\(P_3>P_2=P_1\)}
    {\(P_3=P_2>P_1\)}
    {\(P_3=P_2=P_1\)}
\end{problem}

\begin{answer}
D
\end{answer}

\begin{solution}
设房屋水平投影的面积为 \(S\). 一个斜面与水平面所成的角为 \(\alpha\) 时，该斜面面积乘以 \(\cos\alpha\)，等于它在水平面上的投影面积. 三种盖法的各个斜面投影均覆盖同一个房屋水平投影，且内部没有重叠，因此
\[
P_1\cos\alpha=P_2\cos\alpha=P_3\cos\alpha=S
\]
由于斜面不是竖直面，\(\cos\alpha>0\)，从而 \(P_1=P_2=P_3\). 故选 D.
\end{solution}

\begin{problem}
如图，小圆圈表示网络的结点，结点之间的连线表示它们有网线相联. 连线标注的数字表示该段网线单位时间内可以通过的最大信息量. 现从结点 \(A\) 向结点 \(B\) 传递信息，信息可以分开沿不同的路线同时传递. 则单位时间内传递的最大信息量是（\quad）

\bitmapfigure[max width=0.62\linewidth]{img/2001/new_curriculum_science/q12_fig1.png}

\choices
    {\(26\)}
    {\(24\)}
    {\(20\)}
    {\(19\)}
\end{problem}

\begin{answer}
D
\end{answer}

\begin{solution}
把信息传递方向看作从 \(B\) 到 \(A\)，各网线容量就是图中标注的数值. 先构造总量为 \(19\) 的传输方案：从 \(B\) 出发的四条路径分别传输 \(3\)，\(4\)，\(6\)，\(6\) 个单位. 上方两条路径在右上结点汇合，所占用的容量为 \(3+4=7\)，不超过通往 \(A\) 的容量 \(12\)；下方两条路径在右下结点汇合，所占用的容量为 \(6+6=12\)，恰好不超过通往 \(A\) 的容量 \(12\). 所以最大传输量至少为 \(19\).

另一方面，取一个割：源侧包含 \(B\)、下方的两个中间结点以及右下结点，汇侧包含上方的两个中间结点、右上结点和 \(A\). 该割只穿过容量为 \(3\)，\(4\)，\(12\) 的三条网线，割容量为
\[
3+4+12=19
\]
任何从 \(B\) 到 \(A\) 的传输都必须穿过这组网线，所以最大传输量不超过 \(19\). 上下界相等，故最大信息量为 \(19\)，选 D.
\end{solution}

\section{填空题}
\begin{problem}
若复数 \(z=\sqrt{2}+\sqrt{6}\,\i\)，则 \(\arg\dfrac{1}{z}\) 等于\fillinblank{}.
\end{problem}

\begin{answer}
\(\frac{5\pi}{3}\)
\end{answer}

\begin{solution}
复数 \(z\) 的模为
\[
\lvert z\rvert=\sqrt{(\sqrt2)^2+(\sqrt6)^2}=\sqrt8=2\sqrt2
\]
它的实部、虚部均为正，且
\[
\tan\arg z=\frac{\sqrt6}{\sqrt2}=\sqrt3
\]
所以在 \([0,2\pi)\) 内 \(\arg z=\dfrac\pi3\). 又
\[
\frac1z=\frac{\overline z}{\lvert z\rvert^2}=\frac{\sqrt2-\sqrt6\,\i}{8}
\]
其辐角为 \(-\dfrac\pi3\)，在 \([0,2\pi)\) 内表示为
\[
\arg\frac1z=2\pi-\frac\pi3=\frac{5\pi}{3}
\]
\end{solution}

\begin{problem}
一个袋子里装有大小相同的 \(3\text{ 个}\) 红球和 \(2\text{ 个}\) 黄球. 从中同时取出 \(2\text{ 个}\)，则其中含红球个数的数学期望是\fillinblank{}.
\end{problem}

\begin{answer}
\(1.2\)
\end{answer}

\begin{solution}
设取出的红球个数为随机变量 \(X\)，并将三个红球分别记为 \(R_1\)，\(R_2\)，\(R_3\). 对任意一个红球 \(R_i\)，它被取出的概率为
\[
P(R_i\text{ 被取出})=\frac{\mathrm C_4^1}{\mathrm C_5^2}=\frac25
\]
因为固定取出 \(R_i\) 后，另一个球有 \(4\) 种选择，而同时取出两个球共有 \(\mathrm C_5^2\) 种等可能结果. 令 \(I_i\) 表示红球 \(R_i\) 被取出，取值为 \(0\) 或 \(1\)，则 \(X=I_1+I_2+I_3\). 因此
\[
E(X)=E(I_1)+E(I_2)+E(I_3)=3\times\frac25=\frac65=1.2
\]
所以填 \(1.2\).
\end{solution}

\begin{problem}
在空间中，

\begin{circlelist}
\item 若四点不共面，则这四点中任何三点都不共线.
\item 若两条直线没有公共点，则这两条直线是异面直线.
\end{circlelist}

以上两个命题中，逆命题为真命题的是\fillinblank{}.
\end{problem}

\begin{answer}
\(\circled{2}\)
\end{answer}

\begin{solution}
命题 \(1\) 的逆命题是“若四点中任何三点都不共线，则这四点不共面”. 该逆命题不成立，例如一个平面内的矩形四个顶点中任何三点都不共线，但四点仍然共面.

命题 \(2\) 的逆命题是“若两条直线是异面直线，则这两条直线没有公共点”. 异面直线的定义就是两条既不相交又不平行、且不在同一平面内的直线，特别地它们没有公共点，所以该逆命题是真命题. 因此填 \(\circled{2}\).
\end{solution}

\begin{problem}
设 \(\{a_n\}\) 是公比为 \(q\) 的等比数列，\(S_n\) 是它的前 \(n\) 项和. 若 \(\lvert S_n\rvert\) 是等差数列，则 \(q=\fillinblank{}\).
\end{problem}

\begin{answer}
1
\end{answer}

\begin{solution}
按等比数列的通常约定，首项 \(a_1\ne0\)，公比 \(q\ne0\). 记 \(b_n=\lvert S_n\rvert\). 因为 \(b_n\geqslant0\) 对一切正整数 \(n\) 成立，若 \(\{b_n\}\) 为等差数列，其公差不能小于 \(0\).

当 \(q=1\) 时，\(S_n=na_1\)，从而 \(b_n=n\lvert a_1\rvert\)，确为等差数列.

下面说明其他非零的 \(q\) 均不可能. 当 \(q\ne1\) 时，令 \(C=\left\lvert\dfrac{a_1}{1-q}\right\rvert>0\)，则
\[
b_n=C\lvert1-q^n\rvert
\]
若 \(0<q<1\)，则 \(b_n=C(1-q^n)\)，相邻两项之差为
\[
b_{n+1}-b_n=C(1-q)q^n
\]
该差随 \(n\) 变化，不是常数. 若 \(q>1\)，则 \(b_n=C(q^n-1)\)，相邻两项之差为 \(C(q-1)q^n\)，同样不是常数.

若 \(-1<q<0\)，写 \(q=-r\)，其中 \(0<r<1\). 此时 \(b_1=C(1+r)\)，\(b_2=C(1-r^2)<b_1\)，等差数列公差小于 \(0\)，但非负等差数列不可能具有负公差. 若 \(q=-1\)，则 \(b_1=2C\)，\(b_2=0\)，\(b_3=2C\)，相邻两项之差不相等.

最后，若 \(q<-1\)，写 \(q=-r\)，其中 \(r>1\). 对正整数 \(n\)，
\[
b_{2n+1}-b_{2n}=C\bigl(1+r^{2n+1}-(r^{2n}-1)\bigr)=C\bigl(2+r^{2n}(r-1)\bigr)
\]
该差随 \(n\) 变化，也不可能为等差数列的公差. 综上，只有 \(q=1\) 符合条件.
\end{solution}

\section{解答题}
\begin{problem}
解关于 \(x\) 的不等式 \(\dfrac{x-a}{x-a^2}<0\)（\(a\in\R\)）.
\end{problem}

\begin{solution}
先注意定义域要求 \(x\ne a^2\). 分式小于零，当且仅当分子、分母异号，即原不等式等价于下列两个不等式组的解集的并集：
\[
\begin{cases}
x-a>0,\\
x-a^2<0,
\end{cases}
\qquad\text{或}\qquad
\begin{cases}
x-a<0,\\
x-a^2>0.
\end{cases}
\]

当 \(a<0\) 或 \(a>1\) 时，\(a<a^2\). 此时第一个不等式组的解集为 \(a<x<a^2\)，第二个不等式组无解.

当 \(0<a<1\) 时，\(a^2<a\). 此时第一个不等式组无解，第二个不等式组的解集为 \(a^2<x<a\).

当 \(a=0\) 时，原式为 \(\dfrac{x}{x}=1\)，其定义域为 \(x\ne0\)，不满足小于零；当 \(a=1\) 时，原式为 \(\dfrac{x-1}{x-1}=1\)，其定义域为 \(x\ne1\)，同样无解.

因此不等式的解集为
\[
\begin{cases}
\{x\mid a<x<a^2\},&a<0\text{ 或 }a>1,\\
\{x\mid a^2<x<a\},&0<a<1,\\
\varnothing,&a=0\text{ 或 }a=1.
\end{cases}
\]
\end{solution}

\begin{problem}
如图，用 \(A\)，\(B\)，\(C\) 三类不同的元件连接成两个系统 \(N_1\)，\(N_2\).

\begin{enumerate}
\item 当元件 \(A\)，\(B\)，\(C\) 都正常工作时，系统 \(N_1\) 正常工作；
\item 当元件 \(A\) 正常工作且元件 \(B\)，\(C\) 至少有一个正常工作时，系统 \(N_2\) 正常工作.
\end{enumerate}

已知元件 \(A\)，\(B\)，\(C\) 正常工作的概率依次为 \(0.80\)，\(0.90\)，\(0.90\). 分别求系统 \(N_1\)，\(N_2\) 正常工作的概率 \(P_1\)，\(P_2\).

\begingroup
\leftskip=0pt\relax
\rightskip=0pt\relax
\bitmapfigure[max width=0.74\linewidth]{img/2001/new_curriculum_science/q18_fig1.png}
\endgroup
\end{problem}

\begin{solution}
记元件 \(A\)，\(B\)，\(C\) 正常工作分别为事件 \(A\)，\(B\)，\(C\). 由题设，三个事件相互独立，且
\(P(A)=0.80\)，\(P(B)=0.90\)，\(P(C)=0.90\).

系统 \(N_1\) 正常工作的事件为 \(A\cap B\cap C\)，所以
\[
P_1=P(A\cap B\cap C)=P(A)P(B)P(C)=0.80\times0.90\times0.90=0.648
\]

系统 \(N_2\) 正常工作的事件为 \(A\cap(B\cup C)\). 其对立事件为 \(\overline A\cup(\overline B\cap\overline C)\)，直接按独立性计算更方便：
\[
P_2=P(A)P(B\cup C)=P(A)\bigl[1-P(\overline B\cap\overline C)\bigr]
\]
又有
\(P(\overline B)=1-0.90=0.10\)，\(P(\overline C)=1-0.90=0.10\)
所以
\[
P_2=0.80\bigl[1-0.10\times0.10\bigr]=0.80\times0.99=0.792
\]
因此 \(P_1=0.648\)，\(P_2=0.792\).
\end{solution}

\begin{problem}
设 \(a>0\)，\(f(x)=\dfrac{\e^x}{a}+\dfrac{a}{\e^x}\) 是 \(\R\) 上的偶函数.

\begin{enumerate}
\item 求 \(a\) 的值；
\item 证明 \(f(x)\) 在 \((0,+\infty)\) 上是增函数.
\end{enumerate}
\end{problem}

\begin{solution}
\begin{enumerate}
\item 因为 \(f(x)\) 是偶函数，所以对一切 \(x\in\R\)，有 \(f(x)=f(-x)\). 代入函数表达式得
\[
\frac{\e^x}{a}+\frac{a}{\e^x}=\frac{\e^{-x}}{a}+a\e^x
\]
移项并整理，得到
\[
\left(\frac1a-a\right)\left(\e^x-\e^{-x}\right)=0
\]
取 \(x=1\)，由于 \(\e-\e^{-1}\ne0\)，所以 \(\dfrac1a-a=0\)，即 \(a^2=1\). 又因为 \(a>0\)，故 \(a=1\).

\item 由上问，\(f(x)=\e^x+\e^{-x}\). 在 \((0,+\infty)\) 上，
\[
f'(x)=\e^x-\e^{-x}=\e^{-x}\left(\e^{2x}-1\right)
\]
当 \(x>0\) 时，\(\e^{-x}>0\)，且 \(\e^{2x}>1\)，所以 \(f'(x)>0\). 因此 \(f(x)\) 在 \((0,+\infty)\) 上是增函数.
\end{enumerate}
\end{solution}

\section{选做题：解答题}

\begin{problem}[甲]
如图，以正四棱锥 \(V-ABCD\) 底面中心 \(O\) 为坐标原点建立空间直角坐标系 \(O-xyz\)，其中 \(Ox\parallel BC\)，\(Oy\parallel AB\). \(E\) 为 \(VC\) 中点，正四棱锥底面边长为 \(2a\)，高为 \(h_0\).

\begin{enumerate}
\item 求 \(\cos\langle\overrightarrow{BE},\overrightarrow{DE}\rangle\)；
\item 记面 \(BCV\) 为 \(\alpha\)，面 \(DCV\) 为 \(\beta\)，若 \(\angle BED\) 是二面角 \(\alpha-VC-\beta\) 的平面角，求 \(\angle BED\).
\end{enumerate}

\bitmapfigure[max width=0.46\linewidth]{img/2001/new_curriculum_science/q20_fig1.png}
\end{problem}

\begin{solution}
由坐标系方向和正方形底面可取
\(B(a,a,0)\)，\(C(-a,a,0)\)，\(D(-a,-a,0)\)，\(V(0,0,h_0)\).
因为 \(E\) 是 \(VC\) 的中点，所以
\[
E=\left(-\frac a2,\frac a2,\frac{h_0}{2}\right)
\]
\begin{enumerate}
\item 由坐标作差，得
\(\overrightarrow{BE}=\left(-\frac{3a}{2},-\frac a2,\frac{h_0}{2}\right)\)，\(\overrightarrow{DE}=\left(\frac a2,\frac{3a}{2},\frac{h_0}{2}\right)\).
因此
\[
\overrightarrow{BE}\cdot\overrightarrow{DE}
=-\frac{3a^2}{4}-\frac{3a^2}{4}+\frac{h_0^2}{4}
=\frac{h_0^2-6a^2}{4}
\]
又有
\[
\lvert\overrightarrow{BE}\rvert^2=\lvert\overrightarrow{DE}\rvert^2
=\frac{9a^2}{4}+\frac{a^2}{4}+\frac{h_0^2}{4}
=\frac{10a^2+h_0^2}{4}
\]
由数量积公式，
\[
\cos\langle\overrightarrow{BE},\overrightarrow{DE}\rangle
=\frac{\overrightarrow{BE}\cdot\overrightarrow{DE}}{\lvert\overrightarrow{BE}\rvert\lvert\overrightarrow{DE}\rvert}
=\frac{h_0^2-6a^2}{10a^2+h_0^2}
\]

\item 二面角的平面角所在的截面垂直于棱 \(VC\)，所以 \(BE\perp VC\)，同时 \(DE\perp VC\). 由
\[
\overrightarrow{CV}=(a,-a,h_0)
\]
和上面得到的 \(\overrightarrow{BE}\)，有
\[
\overrightarrow{BE}\cdot\overrightarrow{CV}
=-\frac{3a^2}{2}+\frac{a^2}{2}+\frac{h_0^2}{2}=0
\]
从而 \(h_0^2=2a^2\)，即 \(h_0=\sqrt2a\). 代入第（1）问的结果，得
\[
\cos\angle BED
=\frac{2a^2-6a^2}{10a^2+2a^2}=-\frac13
\]
因此
\[
\angle BED=\arccos\left(-\frac13\right)=\pi-\arccos\frac13
\]
\end{enumerate}
\end{solution}

\reuseproblemnumber
\begin{problem}[乙]
如图，在底面是直角梯形的四棱锥 \(S-ABCD\) 中，\(\angle ABC=90^\circ\)，\(SA\perp\text{面 }ABCD\)，\(SA=AB=BC=1\)，\(AD=\frac12\).

\begin{enumerate}
\item 求四棱锥 \(S-ABCD\) 的体积；
\item 求面 \(SCD\) 与面 \(SBA\) 所成的二面角的正切值.
\end{enumerate}

\bitmapfigure[max width=0.50\linewidth]{img/2001/new_curriculum_science/q20_fig2.png}
\end{problem}

\begin{solution}
\begin{enumerate}
\item 因为 \(AD\parallel BC\)，且 \(AB\perp BC\)，所以 \(AB\) 是直角梯形的高. 底面面积为
\[
S_{ABCD}=\frac12(AD+BC)\cdot AB=\frac12\left(\frac12+1\right)\cdot1=\frac34
\]
四棱锥的高为 \(SA=1\)，所以体积为
\[
V=\frac13S_{ABCD}\cdot SA=\frac13\cdot\frac34\cdot1=\frac14
\]

\item 延长 \(BA\)，\(CD\)，令二者交于 \(E\). 因为 \(AD\parallel BC\)，三角形 \(EAD\) 与三角形 \(EBC\) 相似，且 \(\dfrac{EA}{EB}=\dfrac{AD}{BC}=\dfrac12\). 又 \(E\)，\(A\)，\(B\) 共线，故 \(EB=EA+AB\)，从而 \(EA=AB=1\). 因此 \(EA=AB=SA=1\). 又 \(SA\perp\text{面 }ABCD\)，故 \(SA\perp EB\)，而 \(SA\perp AB\)，从而在直角三角形 \(SEA\) 中，\(SE\perp SB\).

面 \(SBA\) 与面 \(SCD\) 的交线是 \(SE\)，因为 \(A\)，\(B\)，\(E\) 共线且 \(C\)，\(D\)，\(E\) 共线. 由于 \(SA\perp\text{面 }ABCD\)，而 \(EBC\) 是底面，故面 \(SEB\perp\text{面 }EBC\). 又 \(EB\perp BC\)，且 \(EB\) 是两面的交线，所以 \(BC\perp\text{面 }SEB\). 因为 \(BC\perp\text{面 }SEB\)，点 \(C\) 在面 \(SEB\) 上的正投影为 \(B\)，所以 \(SB\) 是斜线 \(SC\) 在该平面上的射影. 又 \(SB\perp SE\)，且 \(SE\) 在面 \(SEB\) 内，根据三垂线定理可得 \(SC\perp SE\). 因此所求二面角的平面角为 \(\angle BSC\).

在直角三角形 \(SBC\) 中，
\(SB=\sqrt{SA^2+AB^2}=\sqrt2\)，\(BC=1\)
且 \(SB\perp BC\). 所以
\[
\tan\angle BSC=\frac{BC}{SB}=\frac1{\sqrt2}=\frac{\sqrt2}{2}
\]
\end{enumerate}
\end{solution}

\begin{problem}
某电厂冷却塔的外形如图所示双曲线的一部分绕其中轴（即双曲线的虚轴）旋转所形成的曲面，其中 \(A\)，\(A'\) 是双曲线的顶点，\(C\)，\(C'\) 是冷却塔上口直径的两个端点，\(B\)，\(B'\) 是下底直径的两个端点，已知 \(AA'=14\mathrm{~m}\)，\(CC'=18\mathrm{~m}\)，\(BB'=22\mathrm{~m}\)，塔高 \(20\mathrm{~m}\).

\begin{enumerate}
\item 建立坐标系并写出该双曲线方程；
\item 求冷却塔的容积（精确到 \(10\mathrm{~m^3}\)，塔壁厚度不计，\(\pi\) 取 \(3.14\)）.
\end{enumerate}

\bitmapfigure[max width=0.48\linewidth]{img/2001/new_curriculum_science/q21_fig1.png}
\end{problem}

\begin{solution}
\begin{enumerate}
\item 建立直角坐标系 \(xOy\)，使 \(AA'\) 在 \(x\) 轴上，\(AA'\) 的中点为原点 \(O\)，且 \(CC'\)，\(BB'\) 平行于 \(x\) 轴. 设双曲线方程为
\(\frac{x^2}{a^2}-\frac{y^2}{b^2}=1\)，\((a>0,b>0)\).
由 \(AA'=14\)，得 \(a=7\). 取右支上的点 \(B(11,y_1)\)，\(C(9,y_2)\)，则由直径长度和塔高可知 \(y_2-y_1=20\)，且
\[
\begin{cases}
\dfrac{11^2}{7^2}-\dfrac{y_1^2}{b^2}=1,\\
\dfrac{9^2}{7^2}-\dfrac{y_2^2}{b^2}=1,\\
y_2-y_1=20.
\end{cases}
\]
前两式给出 \(\dfrac{y_1^2}{y_2^2}=\dfrac{72}{32}=\dfrac94\). 因为 \(B\) 在原点下方，\(C\) 在原点上方，故 \(y_1=-\dfrac32y_2\). 联立 \(y_2-y_1=20\)，得
\(y_1=-12\)，\(y_2=8\).
代入第二个双曲线方程，
\[
\frac{81}{49}-\frac{64}{b^2}=1
\]
从而 \(b^2=98\)，即 \(b=7\sqrt2\). 所以双曲线方程为
\[
\frac{x^2}{49}-\frac{y^2}{98}=1
\]

\item 由于 \(CC'\)、\(BB'\) 是垂直于旋转轴的直径，且旋转轴为 \(y\) 轴，所以右支上点 \(C\)、\(B\) 的横坐标分别为 \(x_C=\dfrac{CC'}{2}=9\)、\(x_B=\dfrac{BB'}{2}=11\)，这与上面取 \(C(9,y_2)\)、\(B(11,y_1)\) 一致.

由双曲线方程，得到
\[
x^2=49+\frac12y^2
\]
在高度 \(y\) 处，冷却塔的截面半径为 \(x\)，所以其容积为
\[
V=\pi\int_{-12}^{8}x^2\,\mathrm dy
=\pi\int_{-12}^{8}\left(49+\frac12y^2\right)\,\mathrm dy
\]
计算得
\[
V=\pi\left.\left(49y+\frac16y^3\right)\right|_{-12}^{8}
=\frac{4060\pi}{3}
\]
    取 \(\pi=3.14\)，得 \(V\approx4249.47\mathrm{~m^3}\). 精确到 \(10\mathrm{~m^3}\)，冷却塔容积为
\[
4250\mathrm{~m^3}
\]
\end{enumerate}
\end{solution}

\begin{problem}
设 \(0<\theta<\frac{\pi}{2}\)，曲线 \(x^2\sin\theta+y^2\cos\theta=1\) 和 \(x^2\cos\theta-y^2\sin\theta=1\) 有 \(4\) 个不同的交点.

\begin{enumerate}
\item 求 \(\theta\) 的取值范围；
\item 证明这 \(4\) 个交点共圆，并求圆半径的取值范围.
\end{enumerate}
\end{problem}

\begin{solution}
两曲线的交点 \((x,y)\) 满足
\[
\begin{cases}
x^2\sin\theta+y^2\cos\theta=1,\\
x^2\cos\theta-y^2\sin\theta=1.
\end{cases}
\]
\begin{enumerate}
\item 令 \(X=x^2\)，\(Y=y^2\). 上述方程组的系数矩阵为
\(M=\begin{pmatrix}\sin\theta&\cos\theta\\\cos\theta&-\sin\theta\end{pmatrix}\)，且 \(M^2=I\). 因而
\[
\begin{pmatrix}X\\Y\end{pmatrix}
=M\begin{pmatrix}1\\1\end{pmatrix}
=\begin{pmatrix}\sin\theta+\cos\theta\\\cos\theta-\sin\theta\end{pmatrix}
\]
即
\(X=\sin\theta+\cos\theta\)，\(Y=\cos\theta-\sin\theta\).
有四个不同交点，当且仅当 \(X>0\) 且 \(Y>0\)，因为这时 \(x\) 和 \(y\) 分别有两个非零取值. 在 \(0<\theta<\dfrac\pi2\) 内，\(\sin\theta+\cos\theta>0\) 恒成立，而 \(\cos\theta-\sin\theta>0\) 等价于 \(\theta<\dfrac\pi4\). 所以
\[
0<\theta<\frac\pi4
\]

\item 将已求得的两个等式 \(x^2=\sin\theta+\cos\theta\) 与 \(y^2=\cos\theta-\sin\theta\) 相加，得
\[
x^2+y^2=2\cos\theta
\]
所以四个交点都满足圆的方程
\[
x^2+y^2=2\cos\theta
\]
即它们共圆，圆心为原点，半径为
\[
r=\sqrt{2\cos\theta}
\]
当 \(0<\theta<\dfrac\pi4\) 时，\(\cos\theta\) 严格递减，且
\(\cos0=1\)，\(\cos\frac\pi4=\frac{\sqrt2}{2}\).
由于区间端点不取，得到
\[
\sqrt[4]{2}<r<\sqrt2
\]
\end{enumerate}
\end{solution}
